\documentclass{article}
\usepackage{graphicx} % Required for inserting images
\usepackage[utf8]{inputenc} % или \usepackage{fontspec} для XeLaTeX/LuaLaTeX
\usepackage[russian]{babel}
\usepackage{geometry}
\geometry{left=3.5cm, right=3.5cm, top=3cm, bottom=1cm}
\usepackage{tikz}
\usetikzlibrary{angles}
\usepackage{fancyhdr}
\usepackage{ccicons}
\usepackage{mathpazo} %Шрифт
\usepackage{\xspace}
\usepackage{wrapfig}


\fancyfoot{}
\fancyhead{}
\pagestyle{fancy}
\renewcommand{\headrulewidth}{0pt}
\fancypagestyle{notitle}{\fancyhead[C]{\Large КНИГА I ПРЕДЛ. IV. ТЕОРЕМА}
\fancyhead[L]{\Large 28}}



\begin{document}
\begin{center}
    \fontsize{20}{25}\selectfontПЕРВЫЕ ШЕСТЬ КНИГ\\\fontsize{25}{25}\selectfontНАЧАЛ ЕВКЛИДА\\
    \vspace{0.4cm}
    \largeВ КОТОРЫХ ИСПОЛЬЗУЮТСЯ ЦВЕТНЫЕ СХЕМЫ И ЗНАКИ ВМЕСТО БУКВ ДЛЯ БОЛЬШЕГО УДОБСТВА ОБУЧАЮЩИХСЯ\\
    \vspace{0.5cm}
    \fontsize{20}{}\selectfontОЛИВЕРА БИРНА\\
    \vspace{2cm}
    
    \begin{tikzpicture}[scale=0.85]
        \draw[fill=black, draw=none] (5, 5) ++(143.13:1cm) arc(143.13:180:1cm) -- (5,5) -- cycle
        \draw[red, line width=7pt] (0,5) -- (5,5)
        \draw[yellow, line width=7pt] (0,5) -- (1.8,7.4)
        \draw[blue, line width = 7pt] (5,5) -- (1.8,7.4)
        
        \fill[orange] (0,0) rectangle (5,5)
        \fill[black, rotate around ={53:(0,5)}] (0, 5) rectangle (3, 8)
        \fill[red, rotate around = {53:(5,5)}] (5, 5) rectangle (9, 9)
        \fill[blue] (0,0) rectangle (1.8, 5)
        
        \draw[white, line width = 3] (1.8, 0) -- (1.8, 7.4)
        \draw[red, dash pattern=on 4pt off 2pt, line width = 3] (1.8, 0) -- (1.8, 7.4)
        \draw[white, line width = 3] (5, 0) -- (5, 5)
        \draw[black, dash pattern=on 4pt off 2pt, line width = 3] (5, 0) -- (5, 5)
        \draw[white, line width = 3] (5, 5) -- (7.4, 8.2)
        \draw[blue, dash pattern=on 4pt off 2pt, line width = 3] (5, 5) -- (7.4, 8.2)
        \draw[black, line width = 3pt] (5,0) -- (1.8, 7.4)
        \draw[black, line width = 3pt] (0, 5) -- (7.4, 8.2)
        \draw[fill=blue, draw=none] (5,5) ++(180:1cm) arc(180:270:1cm) -- (5,5) -- cycle
        \draw[fill=yellow, draw=none] (5,5) ++(53.13:1cm) arc(53.13:143.13:1cm) -- (5,5) -- cycle
    \end{tikzpicture} 
\end{center}

\begin{center}
    \vspace{2.5cm}
    \Large{\fontfamily{ppl}\selectfont
    \textbf{github.com/jemmybutton\\
    2018 изд. 0.3}\\
    }
    \vspace{0.5cm}
    \ccLogo\ \ccAttribution\ \ccShareAlike\\
    \Large{\fontfamily{ppl}\selectfont
    Эта редакция и перевод книги Оливера Бирна The first six
    books of the Elements of Euclid подготовлены Сергеем
    Слюсаревым и распространяются под лицензией CC-BY-SA 4.0
    }
\end{center}

\newpage
\pagestyle{notitle}
\newgeometry{left=7cm, right=2cm, top=3cm, bottom=1cm}

\Large
\begin{wrapfigure}[3]{L}{0.15\textwidth}
\raisebox{10\baselineskip}{\includegraphics[height=3\baselineskip]{E.png}}
\end{wrapfigure}

\notide \hspace{-0.5cm}сли два треугольника имеют по две стороны, равные каждая каждой, (\tikz[baseline=-0.5ex]{\draw[red, line width=4pt] (0, 0) node[above, black, font=\normalsize] {$A$} -- (2.5,0) node[above,black, font=\normalsize ] {$B$}} = \tikz[baseline=-0.5ex]{\draw[red, line width=2pt] (0, 0) node[above, black, font=\normalsize] {$D$} -- (2.5,0) node[above,black, font=\normalsize ] {$E$}} и \tikz[baseline=-0.5ex]{\draw[blue, line width=4pt] (0, 0) node[above, black, font=\normalsize] {$C$} -- (2.5,0) node[above,black, font=\normalsize ] {$A$}} = \tikz[baseline=-0.5ex]{\draw[blue, line width=2pt] (0, 0) node[above,black, font=\normalsize] {$F$} -- (2.5,0) node[above=2pt,black, font=\normalsize ] {$D$}})  \и по равному углу (\tikz[baseline=-3ex]{\draw[fill=orange!70, draw=none] (0,0) ++(235:1cm) arc (235:275:1cm) -- (0,00) -- cycle
\node at (235:1cm) [font=\scriptsize, yshift=-5pt] {$A$};  % точка на первой стороне
\node at (0,0) [font=\scriptsize, yshift=6pt] {$B$};  % вершина (немного смещена вниз)
\node at (275:1cm) [font=\scriptsize, yshift=-5pt] {$C$};  % точка на второй стороне
} = \tikz[baseline=-3ex]{\draw[fill=orange!70, draw=none] (0,0) ++(235:1cm) arc (235:275:1cm) -- (0,00) -- cycle
\node at (235:1cm) [font=\scriptsize, yshift=-5pt] {$D$};  % точка на первой стороне
\node at (0,0) [font=\scriptsize, yshift=6pt] {$E$};  % вершина (немного смещена вниз)
\node at (275:1cm) [font=\scriptsize, yshift=-5pt] {$F$};  % точка на второй стороне
}) содержащемуся между равными прямыми, то они будут иметь и основание равное
основанию (\tikz[baseline=-0.5ex]{\draw[black, line width=4pt] (0,0) node[black, above=0.5pt, font=\normalsize] {$B$} -- (2.5,0) node[black, above=0.5pt, font=\normalsize] {$C$} } = \tikz[baseline=-0.5ex]{\draw[black, line width=2.5pt] (0,0) node[black, above=0.5pt, font=\normalsize] {$E$} -- (2.5,0) node[black, above=0.5pt, font=\normalsize] {$F$} }), и один треугольник
будет равен другому, и остальные углы, стягиваемые
равными сторонами, будут равны каждый каждому (\tikz[baseline=1ex]{\draw[fill=blue, draw=none] (0,0) ++(30:1cm) arc (30:60:1cm) -- (0,0) -- cycle
node at (0,0) [font=\scriptsize, yshift=-5pt] {$B$}
node at (30:1cm) [font=\scriptsize, yshift=-5pt] {$C$}
node at (60:1cm) [font=\scriptsize, yshift=5pt] {$A$}
}
=
\tikz[baseline=1ex]{\draw[fill=blue, draw=none] (0,0) ++(30:1cm) arc (30:60:1cm) -- (0,0) -- cycle
node at (0,0) [font=\scriptsize, yshift=-5pt] {$E$}
node at (30:1cm) [font=\scriptsize, yshift=-5pt] {$F$}
node at (60:1cm) [font=\scriptsize, yshift=5pt] {$D$}
} и \tikz[baseline=1ex]{\draw[fill=red, draw=none] (0,0) ++(100:1cm) arc (100:200:1cm) -- (0,0) -- cycle
node at (0,0) [font=\scriptsize, yshift=-5pt] {$B$}
node at (100:1cm) [font=\scriptsize, yshift=5pt] {$C$}
node at (200:1cm) [font=\scriptsize, yshift=-5pt] {$A$}} = \tikz[baseline=1ex]{\draw[fill=red, draw=none] (0,0) ++(100:1cm) arc (100:200:1cm) -- (0,0) -- cycle
node at (0,0) [font=\scriptsize, yshift=-5pt] {$E$}
node at (100:1cm) [font=\scriptsize, yshift=5pt] {$F$}
node at (200:1cm) [font=\scriptsize, yshift=-5pt] {$D$}}).\\


Представим, что два треугольника расположены таким образом, что вершина одного из двух равных углов, \tikz[baseline=-3ex]{\draw[fill=orange!70, draw=none] (0,0) ++(235:1cm) arc (235:275:1cm) -- (0,00) -- cycle
\node at (235:1cm) [font=\scriptsize, yshift=-5pt] {$A$};  % точка на первой стороне
\node at (0,0) [font=\scriptsize, yshift=6pt] {$B$};  % вершина (немного смещена вниз)
\node at (275:1cm) [font=\scriptsize, yshift=-5pt] {$C$};  % точка на второй стороне
}
или
\tikz[baseline=-3ex]{\draw[fill=orange!70, draw=none] (0,0) ++(235:1cm) arc (235:275:1cm) -- (0,00) -- cycle
\node at (235:1cm) [font=\scriptsize, yshift=-5pt] {$D$};  % точка на первой стороне
\node at (0,0) [font=\scriptsize, yshift=6pt] {$E$};  % вершина (немного смещена вниз)
\node at (275:1cm) [font=\scriptsize, yshift=-5pt] {$F$};  % точка на второй стороне
}
, совпадает с вершиной другого, и \tikz[baseline=-0.5ex]{\draw[red, line width=4pt] (0, 0) node[above=2pt, black, font=\normalsize] {$A$} -- (2.5,0) node[above=2pt,black, font=\normalsize ] {$B$}} совпадает с \tikz[baseline=-0.5ex]{\draw[red, line width=2pt] (0, 0) node[above=2pt, black, font=\normalsize] {$D$} -- (2.5,0) node[above=2pt,black, font=\normalsize ] {$E$}}, тогда \tikz[baseline=-0.5ex]{\draw[blue, line width=4pt] (0, 0) node[above=2pt, black, font=\normalsize] {$C$} -- (2.5,0) node[above=2pt,black, font=\normalsize ] {$A$}} при наложении совпадет с \tikz[baseline=-0.5ex]{\draw[blue, line width=2pt] (0, 0) node[above=2pt,black, font=\normalsize] {$F$} -- (2.5,0) node[above=2pt,black, font=\normalsize ] {$D$}}. Следовательно \tikz[baseline=-0.5ex]{\draw[black, line width=4pt] (0,0) node[black, above=0.5pt, font=\normalsize] {$B$} -- (2.5,0) node[black, above=0.5pt, font=\normalsize] {$C$} } совпадает с \tikz[baseline=-0.5ex]{\draw[black, line width=2.5pt] (0,0) node[black, above=0.5pt, font=\normalsize] {$E$} -- (2.5,0) node[black, above=0.5pt, font=\normalsize] {$F$} }, или же две прямые будут содержать пространство, что невозможно
(акс. X), следовательно \tikz[baseline=-0.5ex]{\draw[black, line width=4pt] (0,0) node[black, above=0.5pt, font=\normalsize] {$B$} -- (2.5,0) node[black, above=0.5pt, font=\normalsize] {$C$} } = \tikz[baseline=-0.5ex]{\draw[black, line width=2.5pt] (0,0) node[black, above=0.5pt, font=\normalsize] {$E$} -- (2.5,0) node[black, above=0.5pt, font=\normalsize] {$F$}}, \tikz[baseline=1ex]{\draw[fill=blue, draw=none] (0,0) ++(30:1cm) arc (30:60:1cm) -- (0,0) -- cycle
node at (0,0) [font=\scriptsize, yshift=-5pt] {$B$}
node at (30:1cm) [font=\scriptsize, yshift=-5pt] {$C$}
node at (60:1cm) [font=\scriptsize, yshift=5pt] {$A$}
} = \tikz[baseline=1ex]{\draw[fill=blue, draw=none] (0,0) ++(30:1cm) arc (30:60:1cm) -- (0,0) -- cycle
node at (0,0) [font=\scriptsize, yshift=-5pt] {$E$}
node at (30:1cm) [font=\scriptsize, yshift=-5pt] {$F$}
node at (60:1cm) [font=\scriptsize, yshift=5pt] {$D$}
}
и 
\tikz[baseline=1ex]{\draw[fill=red, draw=none] (0,0) ++(100:1cm) arc (100:200:1cm) -- (0,0) -- cycle
node at (0,0) [font=\scriptsize, yshift=-5pt] {$B$}
node at (100:1cm) [font=\scriptsize, yshift=5pt] {$C$}
node at (200:1cm) [font=\scriptsize, yshift=-5pt] {$A$}}
=
\tikz[baseline=1ex]{
\draw[fill=red, draw=none] (0,0) ++(100:1cm) arc (100:200:1cm) -- (0,0) -- cycle
node at (0,0) [font=\scriptsize, yshift=-5pt] {$E$}
node at (100:1cm) [font=\scriptsize, yshift=5pt] {$F$}
node at (200:1cm) [font=\scriptsize, yshift=-5pt] {$D$}}
, и поскольку треугольники 
\tikz[baseline=5ex]{
\draw[black, line width=3pt]{(0,0) -- (1.5,0.75)}
\draw[red, line width = 3pt]{(0,0) -- (1.3,2)}
\draw[blue, line width = 3pt] {(1.5,0.7) -- (1.3,2)}
\node at (0,0) [font=\small, black, yshift=-6pt] {$B$}
\node at (1.5,0.75) [font=\small, black, yshift=-7pt] {$C$}
\node at (1.3,2) [font=\small, black, yshift=5pt] {$A$}
}
и
\tikz[baseline=5ex]{
\draw[black, line width=2pt]{(0,0) -- (1.5,0.75)}
\draw[red, line width = 2pt]{(0,0) -- (1.3,2)}
\draw[blue, line width = 2pt] {(1.5,0.7) -- (1.3,2)}
\node at (0,0) [font=\small, black, yshift=-6pt] {$E$}
\node at (1.5,0.75) [font=\small, black, yshift=-7pt] {$F$}
\node at (1.3,2) [font=\small, black, yshift=5pt] {$D$}
}
совпадают при наложении, они
равны во всех отношениях.
ч. т. д.

%Треугольники на полях

%Первый треугольник на полях
\begin{tikzpicture}[remember picture, overlay]
\draw[fill=blue, draw=none] (-6.5,15) ++(25:1cm) arc (25:52:1cm) -- (-6.5,15) -- cycle
\draw[fill=orange!70, draw=none] (-3,20) ++(235:1cm) arc (235:278:1cm) -- (-3,20) -- cycle
\draw[fill=red, draw=none] (-2.5,17) ++(100:0.8cm) arc (100:205:0.8cm) -- (-2.5,17) -- cycle

\draw[black, line width = 3pt] (-6.5,15) -- (-2.5,17)
\draw[red, line width = 3pt] (-6.5,15) -- (-3,20)
\draw[blue, line width = 3pt] (-2.5,17) -- (-3,20)

\node at (-6.5,15) [font=\Large, yshift=-7pt] {$B$}
\node at (-1.5,17) [font=\Large, yshift=-7pt] {$C$}
\node at (-3,20) [font=\Large, yshift=7pt] {$A$}

%Второй треугольник
\draw[fill=blue, draw=none] (-6.5,10) ++(25:1cm) arc (25:52:1cm) -- (-6.5,10) -- cycle
\draw[fill=orange!70, draw=none] (-3,15) ++(235:1cm) arc (235:278:1cm) -- (-3,15) -- cycle
\draw[fill=red, draw=none] (-2.5,12) ++(100:0.8cm) arc (100:205:0.8cm) -- (-2.5,12) -- cycle

\draw[black, line width = 2pt] (-6.5,10) -- (-2.5,12)
\draw[red, line width = 2pt] (-6.5,10) -- (-3,15)
\draw[blue, line width = 2pt] (-2.5,12) -- (-3,15)

\node at (-6.5,10) [font=\Large, yshift=-7pt] {$E$}
\node at (-1.5,12) [font=\Large, yshift=-7pt] {$F$}
\node at (-3,15) [font=\Large, yshift=7pt] {$D$}
\end{tikzpicture}


\end{document}
